\documentclass[UTF8]{ctexart}
\everymath{\displaystyle}

\begin{document}
	\noindent 基：$n$个线性无关的向量$\overrightarrow{a_n}$，可按序数与$n$个系数$b_n$相乘 \\
	以表示$n$维空间中的向量 \\
	条件：\begin{enumerate}
		\item 基是线性无关的
		\item 空间中的所有向量，都可由基线性表示
	\end{enumerate}
	满足任意一个即可
	
	\hspace*{\fill}
	
	\noindent 正交基：内积为零 \\
	内积定义：$<a_n, b_n> = \sum_{i=0}^{n-1} a_i b_i$
	
	\hspace*{\fill}
	
	\noindent 求解向量$\overrightarrow{vec}$在正交基$\overrightarrow{base_n}$下的系数 \\
	系数 = $\frac {\overleftarrow{vec}, \overleftarrow{base_n}}{\overleftarrow{base_n}, \overleftarrow{base_n}}$ $反变换vec=\sum 系数_n\overleftarrow{base_n}$
	
	\hspace*{\fill}
	
	周期函数的内积：$<f(t), g(t)> = \displaystyle \int_T f(t)g(t)$，$T$为$f(t)$周期$T_{f(t)}$和$g(t)$周期$T_{g(t)}$的最小公倍数
	
	三角函数组成的正交函数集：$(\cos(0 \cdot \omega_0), \cos(1 \cdot \omega_0 t) \cdots \cos(n \cdot \omega_0 t),$
	
	傅里叶级数：使用一组正交基
	
	
	1、条件
	2、正交基
	3、
	
	
	
	
\end{document}
